% General Academic Paper Skeleton
% Adaptable to any venue — replace placeholders marked with <...>
\documentclass[11pt,a4paper]{article}

\usepackage[margin=1in]{geometry}
\usepackage[utf8]{inputenc}
\usepackage[T1]{fontenc}
\usepackage{hyperref}
\usepackage{url}
\usepackage{booktabs}
\usepackage{amsfonts}
\usepackage{amsmath}
\usepackage{amssymb}
\usepackage{nicefrac}
\usepackage{microtype}
\usepackage{graphicx}
\usepackage{xcolor}
\usepackage{natbib}
\usepackage{algorithm}
\usepackage{algorithmic}

% Custom commands
\newcommand{\ours}{\textsc{<MethodName>}}
\newcommand{\todo}[1]{\textcolor{red}{[TODO: #1]}}
\newcommand{\R}{\mathbb{R}}
\newcommand{\E}{\mathbb{E}}

\title{<Paper Title>}

\author{
  <Author 1>\thanks{Equal contribution.} \\
  <Affiliation 1> \\
  \texttt{<email1>} \\
  \and
  <Author 2> \\
  <Affiliation 2> \\
  \texttt{<email2>} \\
}

\date{}

\begin{document}

\maketitle

\begin{abstract}
% Structure: Problem → Gap → Method → Results → Impact
\todo{Write abstract — 150-250 words}
\end{abstract}

%% ============================================================
\section{Introduction}
\label{sec:intro}

% Funnel structure: broad → specific → gap → our work → contributions
\todo{Opening paragraph — broad context}

\todo{Problem specification — what exactly needs solving}

\todo{Existing approaches and why they fall short}

\todo{Our approach — key idea in one paragraph}

\noindent Our main contributions are:
\begin{enumerate}
    \item \todo{Contribution 1 — conceptual/theoretical}
    \item \todo{Contribution 2 — methodological}
    \item \todo{Contribution 3 — empirical}
\end{enumerate}

%% ============================================================
\section{Related Work}
\label{sec:related}

\subsection{<Research Direction 1>}
\todo{Group of related work with clear differentiation}

\subsection{<Research Direction 2>}
\todo{Another group — end each subsection with how we differ}

%% ============================================================
\section{Preliminaries}
\label{sec:prelim}

\todo{Background knowledge, notation table, problem formulation}

% Notation table example:
% \begin{table}[h]
% \caption{Notation summary.}
% \centering
% \begin{tabular}{cl}
% \toprule
% Symbol & Description \\
% \midrule
% $x$ & Input \\
% $y$ & Output \\
% $\theta$ & Model parameters \\
% \bottomrule
% \end{tabular}
% \end{table}

%% ============================================================
\section{Method}
\label{sec:method}

% Start with overview paragraph + figure reference
\todo{Method overview — ``Figure~\ref{fig:overview} illustrates...''}

% \begin{figure}[t]
%     \centering
%     % \includegraphics[width=\linewidth]{figures/overview.pdf}
%     \caption{Overview of \ours. \todo{Pipeline figure description.}}
%     \label{fig:overview}
% \end{figure}

\subsection{Problem Formulation}
\label{sec:formulation}
\todo{Formal setup: input space, output space, objective}

\subsection{<Core Component 1>}
\label{sec:comp1}
\todo{Detail with equations — justify each design choice}

\subsection{<Core Component 2>}
\label{sec:comp2}
\todo{Detail with equations}

\subsection{Optimization}
\label{sec:optimization}
\todo{Loss function, training procedure}

% Algorithm example:
% \begin{algorithm}[t]
% \caption{Training procedure for \ours}
% \label{alg:training}
% \begin{algorithmic}[1]
% \REQUIRE Dataset $\mathcal{D}$, learning rate $\eta$
% \ENSURE Trained model parameters $\theta^*$
% \STATE Initialize $\theta$
% \FOR{epoch $= 1$ to $T$}
%     \FOR{mini-batch $(x, y) \sim \mathcal{D}$}
%         \STATE Compute loss $\mathcal{L}(\theta; x, y)$
%         \STATE Update $\theta \leftarrow \theta - \eta \nabla_\theta \mathcal{L}$
%     \ENDFOR
% \ENDFOR
% \end{algorithmic}
% \end{algorithm}

%% ============================================================
\section{Theoretical Analysis}
\label{sec:theory}
% Optional — include if you have theoretical contributions

\todo{Theorem statements, proof sketches (defer full proofs to appendix)}

%% ============================================================
\section{Experiments}
\label{sec:experiments}

\subsection{Setup}
\label{sec:setup}

\paragraph{Datasets.}
\todo{Names, sizes, splits, preprocessing}

\paragraph{Baselines.}
\todo{Methods compared against, with citations}

\paragraph{Metrics.}
\todo{Evaluation metrics with definitions}

\paragraph{Implementation Details.}
\todo{Architecture, hyperparameters, hardware, runtime}

\subsection{Main Results}
\label{sec:main_results}

\todo{Primary comparison — table + analysis paragraphs}

% Table template:
% \begin{table}[t]
% \caption{Comparison on <benchmark>. \textbf{Bold}: best. \underline{Underline}: second best.}
% \label{tab:main}
% \centering
% \begin{tabular}{lccc}
% \toprule
% Method & Metric 1 $\uparrow$ & Metric 2 $\uparrow$ & Metric 3 $\downarrow$ \\
% \midrule
% Baseline 1 & 00.0 & 00.0 & 00.0 \\
% Baseline 2 & 00.0 & 00.0 & 00.0 \\
% \midrule
% \ours & \textbf{00.0} & \textbf{00.0} & \textbf{00.0} \\
% \bottomrule
% \end{tabular}
% \end{table}

\subsection{Ablation Study}
\label{sec:ablation}
\todo{Systematic removal/variation of components}

\subsection{Analysis}
\label{sec:analysis}
\todo{Qualitative examples, visualizations, failure cases, scaling}

%% ============================================================
\section{Conclusion}
\label{sec:conclusion}

% P1: What we did and key takeaways
\todo{Summary of contributions and findings}

% P2: Limitations (be honest)
\todo{Known limitations}

% P3: Future directions
\todo{Future work}

%% ============================================================
\section*{Acknowledgments}
\todo{Funding, compute, helpful discussions}

\bibliographystyle{plainnat}
\bibliography{references}

%% ============================================================
\appendix

\section{Proof of Theorem~\ref{thm:main}}
\label{app:proof}
\todo{Full proof}

\section{Extended Experimental Results}
\label{app:results}
\todo{Additional tables, per-category breakdowns}

\section{Hyperparameter Sensitivity}
\label{app:hyperparams}
\todo{Sensitivity analysis}

\section{Broader Impact Statement}
\label{app:impact}
\todo{Societal implications}

\end{document}
